\documentclass[border=10pt]{standalone}
\usepackage{tikz}
\usepackage{circuitikz}
\usetikzlibrary{positioning, calc, decorations.markings}

% Datapath with ONE shared adder: A = B + C  and  A = A + 3.
% Two 2-to-1 multiplexers (shared select "Mux") choose the adder operands.
% Same building blocks and styles as lect_9_datapath_combine / lect_9_datapath_A_B_plus_C.

\tikzset{register/.style={muxdemux, muxdemux def={Lh=2.4, Rh=2.4, NL=2, NB=1, NT=1,w=4.8,NR=0},
                            muxdemux label={L1=Load, L2=Clock, R1=right1, R2=right,T1=$D_{7..0}$,B1=$Q_{7..0}$,},},
                 }

\tikzset{adder/.style={muxdemux, muxdemux def={Lh=1, Rh=1,NL=0, NB=1, NT=2,w=2,NR=0}}}

\begin{document}
\makeatletter\def\pgf@circ@font@tiny{\footnotesize}\makeatother
\begin{circuitikz}[
    font=\sffamily,
    arrow/.style={-Latex, line width=0.5pt},
    label_text/.style={font=\small},
    bus/.style={-Latex, line width=1.5pt,
        postaction={decorate},
        decoration={markings, mark=at position 0.5 with {
            \draw[thick,-] (-2pt,-3pt) -- (2pt,3pt);
            \node[above left=0.5pt] {\small $8$};
        }}},
    bus late/.style={-Latex, line width=1.5pt, postaction={decorate}, decoration={markings, mark=at position 0.7 with {
            \draw[thick,-] (-2pt,-3pt) -- (2pt,3pt);
            \node[above left=0.5pt] {\small $8$};
        }}},
    bus plain/.style={-Latex, line width=1.5pt},
    control_signal/.style={-Latex, line width=0.5pt},
]

    % ---------------------------------------------------------------- components
    \node[register] at (-5.2,6) (A) {$A$};
    \node[register] at ( 0.0,6) (B) {$B$};
    \node[register] at ( 5.2,6) (C) {$C$};

    % multiplexer 1: input 1 = A (left), input 0 = B (right)
    \draw[thick] (-1.4,3.6) -- (0.4,3.6) -- (-0.2,2.9) -- (-0.8,2.9) -- cycle;
    \node[label_text] at (-1.0,3.35) {$1$};
    \node[label_text] at ( 0.0,3.35) {$0$};

    % multiplexer 2: input 1 = '3' (left), input 0 = C (right)
    \draw[thick] (3.8,3.6) -- (5.6,3.6) -- (5.0,2.9) -- (4.4,2.9) -- cycle;
    \node[label_text] at (4.2,3.35) {$1$};
    \node[label_text] at (5.2,3.35) {$0$};

    % the single adder
    \node[adder, box only, anchor=center] at (2.1,0.6) (Add) {+};

    % ---------------------------------------------------------------- data buses
    % D inputs of B and C
    \draw[bus] ($(B.btpin 1)+(0,1.1)$) -- (B.btpin 1);
    \draw[bus] ($(C.btpin 1)+(0,1.1)$) -- (C.btpin 1);

    % registers -> multiplexers
    \draw[bus] (A.bbpin 1) -- ++(0,-1) -| (-1.0,3.6);
    \draw[bus late] (B.bbpin 1) -- (0,3.6);
    \draw[bus late] (C.bbpin 1) -- (5.2,3.6);

    % constant 3
    \node[above,label_text] at (4.2,4.7) (three) {$'3'$};
    \draw[bus plain] (three) -- (4.2,3.6);

    % multiplexers -> adder
    \draw[bus] (-0.5,2.9) -- ++(0,-1) -| (Add.btpin 1);
    \draw[bus] (4.7,2.9) -- ++(0,-1) -| (Add.btpin 2);

    % adder -> A (feedback loop around the left)
    \draw[bus] (Add.bbpin 1) -- ++(0,-0.8) -| (-8.3,8.2) -| (A.btpin 1);

    % ---------------------------------------------------------------- control signals
    % ALoad; short stubs for BLoad / CLoad
    \draw[control_signal] (-10,0 |- A.blpin 1) node[left, label_text] {$ALoad$} -- (A.blpin 1);
    \draw[control_signal] ($(B.blpin 1)+(-0.6,0)$) -- (B.blpin 1);
    \draw[control_signal] ($(C.blpin 1)+(-0.6,0)$) -- (C.blpin 1);

    % common clock rail
    \draw[control_signal] (-10,4.6) node[left, label_text] {$Clock$} -- (2.7,4.6) |- (C.blpin 2);
    \draw[control_signal] (-7.7,4.6) to[short,*-] (-7.7,0 |- A.blpin 2) -- (A.blpin 2);
    \draw[control_signal] (-2.5,4.6) to[short,*-] (-2.5,0 |- B.blpin 2) -- (B.blpin 2);

    % shared multiplexer select
    \draw[control_signal] (-10,2.5) node[left, label_text] {$Mux$} -- (3.3,2.5) |- (4.271,3.05);
    \draw[control_signal] (-1.8,2.5) to[short,*-] (-1.8,3.05) -- (-0.929,3.05);

\end{circuitikz}
\end{document}
